\section{What must be submitted, and what earns each mark}
\label{sec:deliverables}

\subsection{The five files}

One member uploads all of them through the ``Assignment 3 Submission'' link. The brief deducts marks
for wrong file names, so the names below follow the group-name pattern used for Assignments 1 and 2.

{\small
\begin{longtable}{@{}L{2.5cm}L{5.3cm}L{4.9cm}L{2.3cm}@{}}
\toprule
\textbf{Deliverable} & \textbf{File} & \textbf{Limits and contents} & \textbf{Owner} \\
\midrule
\endhead
Project report & \fp{HCMC_COS30049.1_4_Assignment3_Report.pdf} & 8 pages at most, not counting the cover, contents, bibliography and appendix. The sections of Section~\ref{sec:reportoutline}; the video link on the cover & Cam Dan Bui integrates; every member writes a part \\
Source code & \fp{HCMC_COS30049.1_4_Assignment3_Code.zip} & \fp{frontend/} without \fp{node_modules}; \fp{backend/} with the model code and the models; \fp{README.md} & Le Quang Anh Hoang; Cao Tri Tran for the models \\
Demonstration video & \fp{HCMC_COS30049.1_4_Assignment3_Video.mp4} & 7 minutes at most; MP4 or AVI; everything in Section~\ref{sec:video} & Minh Khang Nguyen records; all speak \\
Video link, if the file is too large for Canvas & \fp{HCMC_COS30049.1_4-video-link.txt} & An unlisted YouTube link that opens in a private window; the same link on the report's cover & Le Quang Anh Hoang \\
Meeting minutes & \fp{HCMC_COS30049.1_4_Assignment3_Meeting_Minutes.pdf} & One meeting for every week since the team formed: Weeks 1 to 12 & Cam Dan Bui \\
Contribution form & \fp{HCMC_COS30049.1_4_Assignment3_Contribution_Form.pdf} & Each member's details and contribution, signed & All \\
\bottomrule
\end{longtable}
}

\subsection{Rubric line by line}
\label{sec:rubricmap}

The middle column is the rubric's own description of full points. The last column names the scope
items (Section~\ref{sec:scopetable}) that earn it.

{\small
\begin{longtable}{@{}L{3.3cm}C{0.7cm}L{5.6cm}L{4.4cm}@{}}
\toprule
\textbf{Criterion} & \textbf{Pts} & \textbf{Full points mean} & \textbf{Earned by} \\
\midrule
\endhead
Report 1: formatting and referencing & 1 & Professional cover with title, group and student details; tidy layout; Harvard references & M-16; Section~\ref{sec:reportoutline} \\
Report 2: architecture and implementation & 2 & Detailed architecture; implementation described with the reasons for each technology; clear diagrams & M-16; the diagrams of Section~\ref{sec:architecture} \\
Report 3: API documentation & 2 & Every endpoint documented with request type, parameters, JSON structure and examples & M-16; generated from the OpenAPI schema (Section~\ref{sec:api}) \\
Code 1: execution & 2 & No errors or warnings; every feature runs; quick start; dependencies managed & M-01, M-14; Section~\ref{sec:readme} \\
Code 2: comments & 1 & Significant sections explained & M-14; Section~\ref{sec:coderules} \\
Code 3: structure & 2 & Clear hierarchy; modular; best practices & M-01, M-05; Section~\ref{sec:layout} \\
Functionality 1: AI model integration & 3 & Seamless; real-time predictions; handles various input types & M-02, M-03, M-04, M-13; S-01, S-05, S-09 \\
Functionality 2: front-end and back-end communication & 4 & Smooth, error-free GET and POST; prompt; full retrieval and submission & M-02, M-05, M-06, M-11 \\
Functionality 3: error handling and validation & 3 & Every input validated on both tiers; graceful errors with clear messages; invalid data never processed & M-09, M-10 \\
Functionality 4: API design & 4 & RESTful; efficient; two or more HTTP methods used properly; documented; reliable & M-01, M-02, M-03, M-12; S-04 adds PUT and DELETE \\
Functionality 5: real-time updates and responsiveness & 4 & Elements and charts update without delay; performance holds across devices & M-11; S-06, S-09 \\
Visualisation 1: diversity and relevance & 3 & Three or more diverse chart types that suit the data and explain the model's output & M-07 (bar, scatter, line); S-02, S-03 \\
Visualisation 2: interactivity & 4 & Zoom, filtering, tooltips, real-time updates; smooth & M-07, M-08; S-06, S-07 \\
Visualisation 3: clarity & 2 & Predictions easy to read; labels, legends, titles and units & M-07; the chart contract (Section~\ref{sec:chartcontract}) \\
Visualisation 4: performance and loading & 2 & Quick to load and smooth with large data & M-07; S-12; Section~\ref{sec:budgets} \\
Usability 1: user interface & 3 & Polished and responsive; intuitive navigation; advanced components such as modals and dynamic forms & M-05, M-06; Section~\ref{sec:components} \\
Usability 2: user experience & 3 & Feedback mechanisms, personalised interactions and smooth transitions & M-10, M-11, M-15; S-04 \\
\bottomrule
\end{longtable}
}

\subsection{The High Distinction list}
\label{sec:hdmap}

The brief asks for three or more chart types, interactive charts and ``at least three advanced
functionalities'', and gives five examples. The plan covers all five, and three of them are met by
Must items alone.

\begin{center}
\small
\begin{tabularx}{\textwidth}{@{}L{4.6cm}Y L{2.4cm}@{}}
\toprule
\textbf{The brief's example} & \textbf{What the Explorer does} & \textbf{Items} \\
\midrule
Enhanced visualisations: three or more chart types with zoom and filtering, in D3.js or Plotly.js & Bar, scatter and line charts in Plotly.js with zoom, pan, a range slider, tooltips, a filter bar and click-through; a dot chart and a grouped bar chart added & M-07, M-08; S-02, S-03, S-07 \\
Real-time updates: live data feeds or auto-refreshing charts & The what-if verdict and curve refresh as the user types; a replay feeds an address's history to the charts; a live lookup streams its progress & M-11; S-06; S-09 \\
Custom user input: advanced forms, full validation, dynamic feedback & The what-if form: amount with a unit switch and a slider, direction, time; rules checked on each keystroke and again on the server & M-06, M-09 \\
Error handling: robust, with clear messages on both ends & One error body, one catalogue of states, a next action on every state & M-10 \\
User-experience extras: export of predictions and charts & CSV and JSON of scored transactions; PNG of every chart; a scenario table & M-12; S-08 \\
\bottomrule
\end{tabularx}
\end{center}

\subsection{Promises from Assignments 1 and 2 that this build must keep}
\label{sec:promises}

{\small
\begin{longtable}{@{}L{5.4cm}L{2.6cm}L{6.0cm}@{}}
\toprule
\textbf{Promise} & \textbf{Where made} & \textbf{How Assignment 3 keeps it} \\
\midrule
\endhead
FR-01 to FR-05: select a supported address; timeline; drilldown; a reason, a comparison basis and the limitation line on every flag; five distinguishable feedback states & A1 Table 2.1 & Views V1 to V3 on the real API; journeys J1 to J3; acceptance tests AT-01 to AT-05 \\
AC-01 to AC-03 and FR-05 pass on every snapshot address and are shown in the video & A1 §1.1 & MS2 runs them on all 28 addresses automatically; the video shows them on one \\
The five planned endpoints, with 400, 404 and 422 & A1 Table 2.7 & E3 to E7 keep their paths, methods and bodies (Section~\ref{sec:api}) \\
The screens move to real data ``without changing screen code'' & A1 §10 & The API client is swapped; screen components change only where a view gains a feature \\
NFR-01: units, source and snapshot date on every chart; the limitation line; none of the five banned words; two walkthrough participants read a flag as ``needs investigation'' & A1 Table 2.5 & The chart contract; the wording test on both tiers; the walkthrough (M-15) \\
NFR-03: keyboard navigation, visible focus, contrast, recovery from mistakes & A1 Table 2.5 & Section~\ref{sec:a11y}; the timeline chart keeps its arrow-key control \\
Charts in Plotly.js; the API called with Axios & A1 WBS dictionary, 1.4.3 and 1.4.4 & D-A3-3 and D-A3-4 \\
The three issues of the heuristic inspection: z-scores may confuse; the unusual-days list hides its overflow; the prototype tools cover content & A1 Table 10.2 & Plain words first, z in the table and tooltip; a ``Show all'' control; the tools removed from the build \\
Optional items go through change control; no live blockchain calls from the application & A1 §2.6, §8.5, WBS 1.4.5 & CR-07 and CR-08 (Appendix~\ref{app:cr}); the Must tier never needs the network \\
The integrated product shows the model evaluation and its known limitations & A1 §9 & View V5, models and limits (M-13) \\
The usability walkthrough, moved to ``the Assignment 3 usability study'' & A2 report, objectives table & M-15, in Week 11, with Assignment 1's four tasks and observer sheet \\
\code{predict.py} is the interface the back end wraps; the Assistant is optional and off by default; Assignment 2's numbers do not change & A2 plan §15 & Section~\ref{sec:model}; S-05; W-06 \\
A post-project review of what changed and what the team learned & A1 §9 & Week 12 minutes (Section~\ref{sec:minutes}) \\
\bottomrule
\end{longtable}
}

\subsection{Rules in the brief that are easy to miss}

\begin{itemize}
  \item \textbf{The video carries the marks.} If it cannot be opened when marking happens, most
    components score zero. It must show desktop, tablet and mobile sizes, every interaction, the
    whole path from input to chart, and the key technical decisions and difficulties.
  \item \textbf{The report documents at least four endpoints} and must describe preprocessing,
    postprocessing, error handling, validation and any change made to the model for the web.
  \item \textbf{The zip holds the model code} that the web application uses, with its
    modifications, and a README that covers set-up, libraries, how to run both tiers and how the
    model is configured.
  \item \textbf{Generative AI}: see Section~\ref{sec:needsteam}, item 1. All use is acknowledged
    in the report.
  \item \textbf{Late work} loses 10\% of the available mark per day and scores zero after seven days.
  \item \textbf{Nothing about the work is posted publicly} on Canvas discussions. The GitHub
    repository for Assignment 2 is public; decision D-A3-1 keeps the Assignment 3 repository
    private until marks are released.
\end{itemize}
